\documentclass[11pt]{article}
\usepackage[margin=1in]{geometry}
\usepackage{mathptmx}
\usepackage{courier}
\usepackage[T1]{fontenc}
\usepackage[utf8]{inputenc}
\usepackage{amsmath,amssymb}
\usepackage{enumitem}
\usepackage{hyperref}
\usepackage{xcolor}
\usepackage{setspace}
\usepackage{titlesec}
\usepackage{booktabs}
\usepackage{tcolorbox}
\usepackage{array}
\hypersetup{colorlinks=true,linkcolor=blue,citecolor=blue,urlcolor=blue}

\pdfstringdefDisableCommands{%
  \def\Lambda{Lambda}%
  \def\Omega{Omega}%
  \def\rho{rho}%
  \def\sigma{sigma}%
  \def\mu{mu}%
  \def\Pi{Pi}%
  \def\alpha{alpha}%
  \def\beta{beta}%
  \def\gamma{gamma}%
  \def\tau{tau}%
  \def\theta{theta}%
  \def\kappa{kappa}%
  \def\lambda{lambda}%
  \def\nu{nu}%
  \def\pi{pi}%
  \def\phi{phi}%
  \def\psi{psi}%
  \def\chi{chi}%
  \def\xi{xi}%
  \def\eta{eta}%
  \def\zeta{zeta}%
  \def\varepsilon{epsilon}%
  \def\varphi{phi}%
  \def\mathrm#1{#1}%
  \def\mathbf#1{#1}%
  \def\mathcal#1{#1}%
  \def\mathbb#1{#1}%
  \def\text#1{#1}%
}

\setstretch{1.12}
\titleformat{\section}{\large\bfseries}{\thesection.}{0.5em}{}
\titleformat{\subsection}{\normalsize\bfseries}{\thesubsection.}{0.5em}{}

\begin{document}

\begin{center}
    {\LARGE \textbf{Overview of the Relational Actualism Four-Paper Suite}}\\[1em]
    {\large Joshua F. Sandeman}\\[0.5em]
    {\small Draft overview for the RA suite}
\end{center}

\vspace{1em}

Relational Actualism (RA) is a discrete foundational programme in which the primary ontology is not a continuum manifold endowed with fields, but a finite actualization graph governed by a structured space of potentia and a Benincasa--Dowker growth law in $d=4$. The four papers collected in this suite develop different sectors of that programme, but they are not independent essays. They are designed to be read as parts of a single research architecture, proceeding from foundations to matter, to gravity and cosmology, and finally to complexity, life, and cognition.

\paragraph{Guiding principle.}
\emph{The aim of RA is not to append discreteness to continuum physics, but to treat discreteness as primary and continuum structure as emergent translation.}

The suite is organized as follows.

\begin{itemize}[leftmargin=1.6em,itemsep=0.6em]
    \item \textbf{Paper I: Kernel and Engine.} 
    Establishes the ontological and dynamical core of the programme: the bimodal ontology, the seven axioms, the sequential growth rule, the BDG kernel, dimensional selection, and the actualization framework. This paper provides the foundational language and formal machinery on which the later papers depend.

    \item \textbf{Paper II: Matter, Forces, and Motifs.} 
    Develops the matter-and-force sector from the topology and renewal structure of the graph. Its central aim is to show how stable particle classes, confinement architecture, and force hierarchy may be constrained by a small set of motif and depth structures in the $d=4$ BDG framework.

    \item \textbf{Paper III: Gravity, Severance, and Cosmology.} 
    Develops the macroscopic sector of RA: gravity as the large-scale description of the actualization graph, causal severance as the replacement for singular behavior, severed-link entropy, the discrete boundary law, and the cosmological programme built around a structurally vanishing cosmological constant.

    \item \textbf{Paper IV: Complexity, Life, and the Causal Firewall.} 
    Extends the framework into complexity, biological organization, agency, and cognition. This is the most exploratory paper in the suite, and is intended as a disciplined extension of the RA architecture rather than as a claim of full closure on these topics.
\end{itemize}

\paragraph{How the suite should be read.}
A central methodological principle runs throughout all four papers: the \emph{discrete statement is primary}; the corresponding continuum formulation, where one exists, is treated as its translation under the appropriate RA dictionary rather than as an independent foundational statement. This principle is especially important in the gravitational and cosmological sectors, where familiar continuum notions such as metric area, curvature, and dark-energy fits are often best understood as effective renderings of underlying discrete facts.

\paragraph{Epistemic discipline.}
The suite contains several distinct kinds of claims, and these should not be conflated. Throughout the papers, results are classified (where appropriate) as:
\begin{itemize}[leftmargin=1.6em,itemsep=0.4em]
    \item \textbf{Lean-verified theorem (LV):} formal theorem proved in Lean~4;
    \item \textbf{Derived result (DR):} analytical consequence of the RA axioms, definitions, and previously established results;
    \item \textbf{Computed / numerical evidence (CV):} supported by simulation, numerical evaluation, or computation, but not yet analytically proved;
    \item \textbf{Phenomenological interpolation (PI):} an empirically successful fit, ansatz, or transition model not yet derived from first principles within RA;
    \item \textbf{Conjecture / open problem (Open):} a structurally motivated claim, theorem target, or proposed interpretation not yet established.
\end{itemize}

This distinction matters. Some of the strongest contributions of the suite are exact and structural; others are numerical or phenomenological; still others remain open research targets. One purpose of the suite is to make those distinctions more explicit, not less.

\paragraph{A note on finitude.}
A further unifying commitment of the suite is \emph{Finitary Actuality}: physically realized universe-states are finite actual graphs. Infinite structures may be used as asymptotic or continuum idealizations, but not as ontologically actual states. This commitment shapes the treatment of severance, entropy, cosmology, and complexity across the suite.

\paragraph{What the suite aims to accomplish.}
The goal of the present programme is not to claim that every sector of physics, cosmology, life, or mind has already been fully derived in final form. Rather, it is to show that a single discrete framework can already organize a wide range of domains that are usually treated separately: particle structure, force hierarchy, gravitation, black-hole entropy, cosmology, complex organization, and the conditions for life and cognition. In some sectors the suite reaches theorem-level closure; in others it offers disciplined conjectures, numerical evidence, or phenomenological bridges. The hope is that the reader will judge the programme not by whether every question has already been closed, but by whether the same underlying architecture is doing nontrivial work across domains that are usually disconnected.

\paragraph{Recommended reading order.}
The papers are designed to be read in order. Paper~I is the natural entry point. Readers primarily interested in particle structure may proceed next to Paper~II; those interested in gravitation and cosmology to Paper~III; and those interested in complexity, life, and cognition to Paper~IV. Papers~II--IV assume the ontology and growth framework of Paper~I, and Papers~III--IV also make selective use of results developed in the earlier parts of the suite.

\end{document}
